% ============================================
% Chapter 8: Test and Validation
% ============================================
% NEW chapter --- Adil writes this

\chapter{Test and Validation}\label{ch:test_validation}

\section{Introduction}

This chapter presents the testing and validation methodology applied to the AEGIS Entropy system. Testing is conducted at three levels: unit tests for individual modules, integration tests for the pipeline, and validation scenarios for the complete system.

\section{Unit Tests}

\subsection{Bridge Pipeline Tests}

The Bridge module includes a comprehensive test suite (\texttt{bridge/test\_pipeline.py}) that validates the offline processing pipeline:

\begin{enumerate}
    \item \textbf{Test 1: Nmap XML parsing} --- Verify that XML files are correctly parsed into feature vectors
    \item \textbf{Test 2: ML model loading} --- Confirm that all three models (RF, XGB, IF) and the scaler load successfully
    \item \textbf{Test 3: Prediction pipeline} --- End-to-end test from Nmap XML to classification output
    \item \textbf{Test 4: Dashboard POST} --- Verify that events are correctly pushed to the Flask dashboard
    \item \textbf{Test 5: Local fallback} --- Confirm that events are written locally when the dashboard is unreachable
\end{enumerate}

The results of all five tests are presented in Table~\ref{tab:unit_tests}.

\begin{table}[H]
\centering
\caption{Bridge pipeline test results}
\label{tab:unit_tests}
\begin{tabular}{lll}
\toprule
\textbf{Test} & \textbf{Status} & \textbf{Notes} \\
\midrule
Nmap XML parsing & PASS & All 10 scans parsed correctly \\
ML model loading & PASS & RF, XGB, IF, Scaler loaded \\
Prediction pipeline & PASS & 5/5 Nmap files classified \\
Dashboard POST & PASS & Alerts received by dashboard \\
Local fallback & PASS & Events written to JSONL \\
\bottomrule
\end{tabular}
\end{table}

\subsection{Model Performance Validation}

The retrained models were validated against the project's success criteria, as shown in Table~\ref{tab:model_validation}:

\begin{table}[H]
\centering
\caption{Model performance on test set (retrained environment)}
\label{tab:model_validation}
\begin{tabular}{lcccc}
\toprule
\textbf{Model} & \textbf{Accuracy} & \textbf{F1-Score} & \textbf{FPR} & \textbf{Status} \\
\midrule
Random Forest & 99.911\% & 99.920\% & 0.114\% & PASS \\
XGBoost & 99.914\% & 99.923\% & 0.114\% & PASS \\
Isolation Forest & 24.627\% & 9.464\% & 53.532\% & Expected (see Ch. 4) \\
\bottomrule
\end{tabular}
\end{table}

\section{Integration Tests}

\subsection{Offline Pipeline Validation}

The offline pipeline was validated by processing all 10 Nmap scans through the Bridge module:

\begin{enumerate}
    \item Parse Nmap XML files into feature vectors
    \item Apply StandardScaler normalization
    \item Classify using RF and XGBoost models
    \item Generate detection events with confidence scores
    \item Push events to the dashboard (or log locally)
\end{enumerate}

\subsection{Deception Module Integration}

The deception modules (core\_deception, network\_mutator, traffic\_shaper) were tested for:
\begin{itemize}
    \item Correct JSONL event logging
    \item HTTP POST to dashboard endpoints
    \item Configuration via \texttt{config.py} (shared paths and thresholds)
\end{itemize}

\section{Validation Scenarios}

\subsection{Scenario 1: Fast SYN Scan}

Simulate a fast Nmap SYN scan (\texttt{-sS -T4}) and verify that:
\begin{itemize}
    \item The ML pipeline classifies the traffic as ATTACK with high confidence ($> 0.95$)
    \item The defense subsystem triggers HIGH severity response
    \item Attacker is redirected to honeypot and blacklisted
\end{itemize}

\subsection{Scenario 2: Slow Stealth Scan}

Simulate a slow scan (\texttt{--scan-delay 500ms --max-rate 5}) and verify that:
\begin{itemize}
    \item Detection occurs within the 90-second sliding window
    \item The system correctly identifies low-and-slow behavior
    \item Appropriate response is triggered despite low packet rate
\end{itemize}

\subsection{Scenario 3: Dashboard Connectivity Loss}

Simulate dashboard unavailability and verify that:
\begin{itemize}
    \item Events are written to local JSONL files
    \item No data is lost during the outage
    \item Events are forwarded to the dashboard when connectivity is restored
\end{itemize}

\section{Conclusion}

The AEGIS Entropy system passes all unit tests (5/5), integration tests, and validation scenarios. The detection pipeline achieves F1-Scores exceeding 99\% on the CICIDS2017 dataset, and the defense subsystem responds within 520ms for full-stack activation. The system is ready for live demonstration and evaluation.
